\documentclass[11pt]{article}
\usepackage[a4paper, top=1in, bottom=1in, left=1in, right=1in]{geometry}
\usepackage{amsmath, amssymb}
\usepackage{mdframed}
\usepackage{enumitem}
\usepackage{multicol}
\usepackage{tikz}
\usetikzlibrary{calc}
\usepackage{array}
\usepackage{fancyhdr}
\usepackage[T1]{fontenc}
\usepackage{booktabs}
\usepackage{xcolor}

\setlength{\headheight}{26pt}
\pagestyle{fancy}
\fancyhf{}
\lhead{\textbf{Quadratic Functions: Notes \& Practice}}
\rhead{Name: \underline{\hspace{4cm}}}
\cfoot{\thepage}

\setlength{\parindent}{0pt}
\setlength{\parskip}{5pt}

% -- Key Idea box --
\newmdenv[
  linewidth=0.8pt,
  innerleftmargin=10pt, innerrightmargin=10pt,
  innertopmargin=6pt,   innerbottommargin=6pt,
  skipabove=8pt, skipbelow=8pt,
  backgroundcolor=gray!8
]{keybox}
\newcommand{\KeyIdea}[1]{\begin{keybox}\textbf{Key Idea.} #1\end{keybox}}

% -- Example box --
\newmdenv[
  linewidth=0.6pt,
  innerleftmargin=10pt, innerrightmargin=10pt,
  innertopmargin=5pt,   innerbottommargin=5pt,
  skipabove=6pt, skipbelow=6pt,
  linecolor=blue!50!black
]{exbox}

% -- workspace --
\newcommand{\ws}[1][3cm]{\vspace{#1}}

\begin{document}

% ==============================================================
\begin{center}
  {\Large\textbf{Quadratic Functions}}\\[4pt]
  {\large Three Forms $\cdot$ Transformations $\cdot$ Graphing $\cdot$ Analysis}
\end{center}
\hrule
\vspace{0.4cm}

% ==============================================================
\section*{Part 1 \quad Three Forms of a Quadratic Function}
% ==============================================================

Every quadratic function can be written in three different but equivalent forms.
Each form highlights different features of the parabola.

\vspace{0.3cm}
\begin{center}
\renewcommand{\arraystretch}{2}
\begin{tabular}{|c|c|c|}
\hline
\textbf{Form} & \textbf{Equation} & \textbf{What it tells you} \\
\hline
Vertex Form   & $y = a(x - h)^2 + k$ & Vertex $= (h,\, k)$; direction; stretch \\
\hline
Standard Form & $y = ax^2 + bx + c$  & $y$-intercept $= (0,c)$; axis of symmetry \\
\hline
Intercept Form & $y = a(x - x_1)(x - x_2)$ & $x$-intercepts (zeros) $= x_1$ and $x_2$ \\
\hline
\end{tabular}
\end{center}

\vspace{0.3cm}
\KeyIdea{In all three forms, the sign of $a$ tells you the direction of the
parabola: $a > 0$ opens \textbf{up}; $a < 0$ opens \textbf{down}.
The value of $|a|$ controls width: $|a|>1$ is narrower than $y=x^2$;
$0 < |a| < 1$ is wider.}

\vspace{0.3cm}
\begin{exbox}
\textbf{Example 1.} Identify the form and key features of each function.

\textbf{(a)} $y = 2(x-3)^2 + 1$ \quad Vertex form. \quad Vertex $= (3,1)$, opens up, narrower than $y=x^2$.

\textbf{(b)} $y = x^2 - 5x + 6$ \quad Standard form. \quad $y$-intercept $= (0,6)$.

\textbf{(c)} $y = -2(x-1)(x+4)$ \quad Intercept form. \quad Zeros at $x=1$ and $x=-4$, opens down.
\end{exbox}

% ==============================================================
\section*{Part 2 \quad Transformations of $f(x) = x^2$}
% ==============================================================

A \textbf{transformation} changes the position, size, or orientation of the parent
parabola $y = f(x) = x^2$. There are four main types.

% ----------------------------------------------------------
\subsection*{1. \; Translation (Shift)}
% ----------------------------------------------------------

A translation moves the parabola without changing its shape.

\begin{center}
\renewcommand{\arraystretch}{1.6}
\begin{tabular}{|l|l|l|}
\hline
\textbf{Transformation} & \textbf{Equation} & \textbf{Effect on graph} \\
\hline
Shift right $h$ units  & $g(x) = f(x - h) = (x-h)^2$  & Vertex moves right \\
Shift left  $h$ units  & $g(x) = f(x + h) = (x+h)^2$  & Vertex moves left \\
Shift up    $k$ units  & $g(x) = f(x) + k = x^2 + k$  & Vertex moves up \\
Shift down  $k$ units  & $g(x) = f(x) - k = x^2 - k$  & Vertex moves down \\
\hline
\end{tabular}
\end{center}

\medskip
\textbf{Important:} Inside the parentheses, the sign is opposite to the direction.
$f(x+2)$ shifts \textit{left} 2; $\; f(x-2)$ shifts \textit{right} 2.

\begin{exbox}
\textbf{Example 2.} Start with $f(x)=x^2$.

$g(x) = (x - 3)^2 + 4$: shift \textbf{right 3}, then \textbf{up 4}.
Vertex moves from $(0,0)$ to $(3,4)$.

$g(x) = (x + 2)^2 - 1$: shift \textbf{left 2}, then \textbf{down 1}.
Vertex moves to $(-2,-1)$.
\end{exbox}

% ----------------------------------------------------------
\subsection*{2. \; Vertical Stretch and Shrink}
% ----------------------------------------------------------

Multiplying $f(x)$ by a positive constant $a$ changes the \textit{width}.

\begin{center}
\renewcommand{\arraystretch}{1.5}
\begin{tabular}{|l|l|}
\hline
$|a| > 1$ \quad (e.g., $y = 3x^2$)          & Vertical \textbf{stretch} --- parabola is \textbf{narrower} \\
$0 < |a| < 1$ \quad (e.g., $y = \tfrac{1}{2}x^2$) & Vertical \textbf{shrink} --- parabola is \textbf{wider} \\
\hline
\end{tabular}
\end{center}

\medskip
Every $y$-value is multiplied by $a$.
The point $(1,1)$ on $y = x^2$ becomes $(1,a)$ on $y = ax^2$.

\begin{exbox}
\textbf{Example 3.}

$y = 2x^2$: vertical stretch by factor 2 --- \textbf{narrower} than $y = x^2$.

$y = \tfrac{1}{2}x^2$: vertical shrink by factor $\tfrac{1}{2}$ --- \textbf{wider} than $y = x^2$.
\end{exbox}

% ----------------------------------------------------------
\subsection*{3. \; Reflection}
% ----------------------------------------------------------

\begin{center}
\renewcommand{\arraystretch}{1.5}
\begin{tabular}{|l|l|l|}
\hline
\textbf{Type} & \textbf{Equation} & \textbf{Effect} \\
\hline
Reflect over $x$-axis & $g(x) = -f(x) = -x^2$    & Flips upside down (opens down) \\
Reflect over $y$-axis & $g(x) = f(-x) = (-x)^2 = x^2$ & For $y=x^2$, no visible change \\
\hline
\end{tabular}
\end{center}

\medskip
For a general function, reflection over the $y$-axis replaces $x$ with $-x$.
For asymmetric quadratics like $y = (x-3)^2$, it matters:
$g(x) = f(-x) = (-x-3)^2 = (x+3)^2$ --- the vertex moves from $(3,0)$ to $(-3,0)$.

\begin{exbox}
\textbf{Example 4.} Start with $y = x^2$.

\textbf{Step 1} Reflect over $x$-axis: $y = -x^2$.

\textbf{Step 2} Vertical stretch by 2: $y = -2x^2$.

\textbf{Step 3} Shift right 5: $y = -2(x-5)^2$.

\textbf{Step 4} Shift up 8: $y = -2(x-5)^2 + 8$. \quad Vertex $(5, 8)$.
\end{exbox}

% ----------------------------------------------------------
\subsection*{4. \; Summary Table --- All Transformations}
% ----------------------------------------------------------

Starting from $y = f(x) = x^2$:

\begin{center}
\renewcommand{\arraystretch}{1.6}
\begin{tabular}{|l|l|}
\hline
\textbf{New function} & \textbf{Transformation} \\
\hline
$g(x) = f(x) + k$     & Shift up $k$ units \\
$g(x) = f(x) - k$     & Shift down $k$ units \\
$g(x) = f(x - h)$     & Shift right $h$ units \\
$g(x) = f(x + h)$     & Shift left $h$ units \\
$g(x) = a\,f(x),\; a>1$    & Vertical stretch (narrower) \\
$g(x) = a\,f(x),\; 0<a<1$  & Vertical shrink (wider) \\
$g(x) = -f(x)$        & Reflection over $x$-axis \\
$g(x) = f(-x)$        & Reflection over $y$-axis \\
\hline
\end{tabular}
\end{center}

% ==============================================================
\section*{Part 3 \quad Writing Equations from a Graph}
% ==============================================================

When you see a graph, pick the form that matches the information given.

\textbf{Strategy A --- Use intercept form} when you can read the $x$-intercepts $x_1, x_2$:
\[
  y = a(x - x_1)(x - x_2)
\]
Then plug in one other known point to solve for $a$.

\textbf{Strategy B --- Use vertex form} when you can read the vertex $(h,k)$:
\[
  y = a(x - h)^2 + k
\]
Then plug in one other known point to solve for $a$.

\begin{exbox}
\textbf{Example 5.} A parabola has zeros at $x = 3$ and $x = 7$, and passes through $(5,\,8)$.

\textbf{Step 1} Write intercept form: $y = a(x-3)(x-7)$.

\textbf{Step 2} Plug in $(5,8)$: \quad $8 = a(5-3)(5-7) = a \cdot 2 \cdot (-2) = -4a$.

\textbf{Step 3} Solve: $a = -2$.

\textbf{Result:} $y = -2(x-3)(x-7)$

\textbf{Standard form:} $y = -2(x^2 - 10x + 21) = -2x^2 + 20x - 42$.
\end{exbox}

% ==============================================================
\section*{Part 4 \quad Converting to Vertex Form}
% ==============================================================

Given $y = ax^2 + bx + c$ in standard form, find the vertex $(h,k)$:

\[
  \boxed{h = -\frac{b}{2a}} \qquad \boxed{k = f(h)}
\]

Then rewrite as $y = a(x-h)^2 + k$.

\begin{exbox}
\textbf{Example 6.} Convert $y = -2x^2 + 20x - 42$ to vertex form.

\textbf{Step 1} Find $h$: $\displaystyle h = -\frac{20}{2(-2)} = -\frac{20}{-4} = 5$.

\textbf{Step 2} Find $k$: $k = f(5) = -2(25) + 20(5) - 42 = -50 + 100 - 42 = 8$.

\textbf{Result:} $y = -2(x-5)^2 + 8$. \quad Vertex $(5,8)$.
\end{exbox}

\begin{exbox}
\textbf{Example 7 (Completing the Square).} Convert $y = 2x^2 - 8x + 8$.

\begin{align*}
y &= 2(x^2 - 4x) + 8 \\
  &= 2(x^2 - 4x + 4 - 4) + 8 \\
  &= 2(x-2)^2 - 8 + 8 \\
  &= 2(x-2)^2
\end{align*}

Vertex $(2, 0)$. The parabola touches the $x$-axis at its vertex.
\end{exbox}

% ==============================================================
\section*{Part 5 \quad Domain, Range, and Zeros}
% ==============================================================

For any quadratic $y = ax^2 + bx + c$:

\begin{center}
\renewcommand{\arraystretch}{1.6}
\begin{tabular}{|l|l|}
\hline
\textbf{Feature}  & \textbf{Rule} \\
\hline
Domain            & Always $(-\infty,\, +\infty)$ \\
Range ($a > 0$, opens up)  & $[\,k,\, +\infty)$ where $k$ is the minimum \\
Range ($a < 0$, opens down) & $(-\infty,\, k\,]$ where $k$ is the maximum \\
Zeros (x-intercepts)  & Set $y = 0$ and solve \\
\hline
\end{tabular}
\end{center}

\begin{exbox}
\textbf{Example 8.} $h(x) = 16x^2 - 4$.

Vertex form: $h(x) = 16(x-0)^2 - 4$, so vertex $= (0, -4)$, opens up.

\textbf{Zeros:} $16x^2 - 4 = 0 \Rightarrow x^2 = \tfrac{1}{4} \Rightarrow x = \pm\tfrac{1}{2}$.

\textbf{Domain:} $(-\infty, +\infty)$. \quad \textbf{Range:} $[-4, +\infty)$.
\end{exbox}

\begin{exbox}
\textbf{Example 9.} $y = -(x+5)(x-1)$.

Zeros: $x = -5$ and $x = 1$.

Vertex $x$-coordinate: midpoint of zeros $= \dfrac{-5+1}{2} = -2$.

$y(-2) = -(-2+5)(-2-1) = -(3)(-3) = 9$. \quad Vertex $= (-2, 9)$. Opens down.

\textbf{Domain:} $(-\infty, +\infty)$. \quad \textbf{Range:} $(-\infty,\, 9\,]$.
\end{exbox}

\newpage
% ==============================================================
\begin{center}
  {\Large\textbf{Practice Problems}}\\[4pt]
  \textit{30 problems covering all five topics above.}
\end{center}
\hrule
\vspace{0.4cm}

% ==============================================================
\subsection*{Section A \quad Three Forms (Problems 1--5)}
% ==============================================================

\textbf{1.} Identify the form of each function and state the indicated feature.

\hspace{1cm}(a) $y = -3(x-4)^2 + 7$ \hspace{1cm} Form: \underline{\hspace{3cm}} \hspace{0.5cm} Vertex: \underline{\hspace{2.5cm}}

\hspace{1cm}(b) $y = x^2 - 7x + 10$ \hspace{1.2cm} Form: \underline{\hspace{3cm}} \hspace{0.5cm} $y$-intercept: \underline{\hspace{1.5cm}}

\hspace{1cm}(c) $y = 5(x+2)(x-6)$ \hspace{1.1cm} Form: \underline{\hspace{3cm}} \hspace{0.5cm} Zeros: \underline{\hspace{2.5cm}}

\ws[2cm]

\textbf{2.} Write the function in \textbf{standard form}: $y = (x-3)(x+5)$.

\ws[3cm]

\textbf{3.} Write the function in \textbf{standard form}: $y = -2(x-1)(x-6)$.

\ws[3cm]

\textbf{4.} Write a \textbf{vertex form} equation for a parabola with vertex $(-2, 5)$ that passes through $(0, 1)$.
Find the value of $a$ first.

\ws[3.5cm]

\textbf{5.} A quadratic has $y$-intercept $(0, 12)$ and zeros at $x = 3$ and $x = 4$.
Write the function in intercept form, find $a$, then convert to standard form.

\ws[3.5cm]

% ==============================================================
\subsection*{Section B \quad Translations (Problems 6--10)}
% ==============================================================

\textbf{6.} Starting with $f(x) = x^2$, write the equation after shifting \textbf{right 4 units} and \textbf{up 3 units}.
State the new vertex.

\ws[2.5cm]

\textbf{7.} Starting with $f(x) = x^2$, write the equation after shifting \textbf{left 5 units} and \textbf{down 2 units}.

\ws[2.5cm]

\textbf{8.} Describe the transformations that take $f(x) = x^2$ to $g(x) = (x + 3)^2 - 7$.
State the vertex.

\ws[2.5cm]

\textbf{9.} The graph of $g(x) = f(x - 2) + 6$ is a translation of $f(x) = x^2$.
\; (a) Describe the translation.\quad (b) Write the equation of $g(x)$.\quad (c) State the vertex.

\ws[3cm]

\textbf{10.} Point $(1, 1)$ is on $y = x^2$. After the translation $g(x) = (x-4)^2 + 2$, what are the
new coordinates of this point?

\ws[2.5cm]

% ==============================================================
\subsection*{Section C \quad Vertical Stretch and Shrink (Problems 11--14)}
% ==============================================================

\textbf{11.} Compare each function to $y = x^2$. State whether the parabola is \textit{narrower},
\textit{wider}, or the \textit{same width}, and by what factor.

\hspace{1cm}(a) $y = 5x^2$ \hspace{2.5cm} (b) $y = \tfrac{1}{4}x^2$
\hspace{2.5cm} (c) $y = 0.8\,x^2$

\ws[2.5cm]

\textbf{12.} Starting with $f(x) = x^2$, write the equation after a \textbf{vertical stretch by factor 3},
then a \textbf{shift up 5 units}.

\ws[2.5cm]

\textbf{13.} Starting with $f(x) = x^2$, write the equation after a \textbf{vertical shrink by $\tfrac{1}{2}$},
then a \textbf{shift right 3 units} and \textbf{down 1 unit}.

\ws[2.5cm]

\textbf{14.} The point $(2, 4)$ is on $y = x^2$. What is the corresponding point on $y = \tfrac{3}{4}x^2$?
What about on $y = 4x^2$?

\ws[2.5cm]

% ==============================================================
\subsection*{Section D \quad Reflections (Problems 15--17)}
% ==============================================================

\textbf{15.} Reflect $f(x) = 2(x - 3)^2 + 1$ over the \textbf{$x$-axis}.
Write the new equation and state the new vertex and direction.

\ws[2.5cm]

\textbf{16.} Reflect $f(x) = (x - 4)^2$ over the \textbf{$y$-axis}.
Write the new equation and state the new vertex.

\ws[2.5cm]

\textbf{17.} Starting with $y = x^2$: first reflect over the $x$-axis, then shift left 3, then up 5.
Write the final equation and identify the vertex.

\ws[3cm]

% ==============================================================
\subsection*{Section E \quad Combined Transformations (Problems 18--22)}
% ==============================================================

\textbf{18.} List ALL transformations applied to $f(x) = x^2$ to produce $g(x) = -4(x + 1)^2 - 3$.
State them in the correct order.

\ws[3.5cm]

\textbf{19.} List ALL transformations applied to $f(x) = x^2$ to produce $p(x) = \tfrac{1}{2}(x - 2)^2 + 4$.

\ws[3cm]

\textbf{20.} Write the equation obtained by applying these steps to $y = x^2$ in order:

\hspace{1cm}(1) Reflect over $x$-axis \hspace{1cm} (2) Stretch vertically by factor 2
\hspace{0.5cm}(3) Shift right 1 unit \hspace{0.5cm}(4) Shift down 4 units

\ws[3cm]

\textbf{21.} The function $g(x) = -\tfrac{1}{2}(x + 3)^2 + 5$ is compared to $f(x) = x^2$.
Describe all differences in direction, width, and position.

\ws[3cm]

\textbf{22.} Write the vertex form equation of a parabola that is:
\textit{reflected over the $x$-axis, vertically stretched by 3, shifted right 2, shifted up 6.}
What is the vertex?

\ws[3cm]

% ==============================================================
\subsection*{Section F \quad Writing Equations from Graphs (Problems 23--26)}
% ==============================================================

\textbf{23.} A parabola opens upward, has $x$-intercepts at $x = -1$ and $x = 5$,
and passes through $(2,\,-9)$.
Write the equation in intercept form, then convert to standard form.

\ws[3.5cm]

\textbf{24.} A parabola has vertex $(3, -2)$ and passes through $(5, 6)$.
Write in vertex form, then standard form.

\ws[3.5cm]

\textbf{25.} A parabola has $x$-intercepts $(-4, 0)$ and $(2, 0)$ and vertex $(-1,\,9)$.
Find $a$, write the intercept form, then expand to standard form.

\ws[3.5cm]

\textbf{26.} From a graph: the parabola opens downward with vertex $(1, 4)$ and
$y$-intercept $(0, 3)$.
Find $a$, write the vertex form, and state the domain and range.

\ws[3.5cm]

% ==============================================================
\subsection*{Section G \quad Converting to Vertex Form (Problems 27--28)}
% ==============================================================

\textbf{27.} Convert each function to vertex form using the vertex formula
$h = -\dfrac{b}{2a}$, $\; k = f(h)$.

\hspace{1cm}(a) $y = x^2 - 6x + 5$

\ws[3cm]

\hspace{1cm}(b) $y = 3x^2 + 12x - 1$

\ws[3cm]

\hspace{1cm}(c) $y = -x^2 + 4x + 12$

\ws[3cm]

\textbf{28.} Convert $y = 2x^2 - 12x + 10$ to vertex form by
\textbf{completing the square} (show all steps).
Then state the vertex.

\ws[4cm]

% ==============================================================
\subsection*{Section H \quad Domain, Range, and Zeros (Problems 29--30)}
% ==============================================================

\textbf{29.} For each function, find: \textbf{(i)} the zeros, \textbf{(ii)} the vertex,
\textbf{(iii)} the domain, \textbf{(iv)} the range.

\hspace{1cm}(a) $f(x) = 2x^2 - 8x + 8$

\ws[4cm]

\hspace{1cm}(b) $y = -(x + 5)(x - 1)$

\ws[4cm]

\hspace{1cm}(c) $h(x) = 16x^2 - 4$

\ws[4cm]

\textbf{30.} A quadratic function opens downward, has zeros at $x = -3$ and $x = 7$,
and a maximum value of 20.

\hspace{1cm}(a) Find the $x$-coordinate of the vertex (midpoint of zeros).

\hspace{1cm}(b) Write the vertex form equation.

\hspace{1cm}(c) Find $a$.

\hspace{1cm}(d) State the range.

\ws[4cm]

% ==============================================================
\newpage
\begin{center}
  {\Large\textbf{Answer Key}}
\end{center}
\hrule
\vspace{0.4cm}
% ==============================================================

\begin{multicols}{2}

\textbf{1.}
(a) Vertex form; vertex $(4,7)$\\
(b) Standard form; $y$-int $(0,10)$\\
(c) Intercept form; zeros $x=-2,\,6$

\medskip
\textbf{2.} $y=(x-3)(x+5)=x^2+2x-15$

\medskip
\textbf{3.} $-2(x-1)(x-6)=-2(x^2-7x+6)$\\
$= -2x^2+14x-12$

\medskip
\textbf{4.} $y=a(x+2)^2+5$. Plug in $(0,1)$:\\
$1=a(4)+5 \Rightarrow a=-1$.\\
$y=-(x+2)^2+5$

\medskip
\textbf{5.} $y=a(x-3)(x-4)$. Plug in $(0,12)$:\\
$12=a(3)(4)=12a \Rightarrow a=1$.\\
$y=(x-3)(x-4)=x^2-7x+12$

\medskip
\textbf{6.} $g(x)=(x-4)^2+3$; vertex $(4,3)$

\medskip
\textbf{7.} $g(x)=(x+5)^2-2$; vertex $(-5,-2)$

\medskip
\textbf{8.} Left 3, down 7; vertex $(-3,-7)$

\medskip
\textbf{9.} (a) Right 2, up 6\\
(b) $g(x)=(x-2)^2+6$\\
(c) Vertex $(2,6)$

\medskip
\textbf{10.} $(1,1)$ shifts to $(1+4,\,1+2)=(5,3)$

\medskip
\textbf{11.}
(a) $a=5>1$: narrower, factor 5\\
(b) $a=\frac{1}{4}<1$: wider, factor $\frac{1}{4}$\\
(c) $a=0.8<1$: wider, factor $0.8$

\medskip
\textbf{12.} $g(x)=3x^2+5$

\medskip
\textbf{13.} $g(x)=\tfrac{1}{2}(x-3)^2-1$

\medskip
\textbf{14.} On $\tfrac{3}{4}x^2$: $(2,\,3)$. On $4x^2$: $(2,16)$

\medskip
\textbf{15.} $g(x)=-2(x-3)^2-1$;\\ vertex $(3,-1)$, opens down

\medskip
\textbf{16.} $g(x)=(-x-4)^2=(x+4)^2$;\\ vertex $(-4,0)$

\medskip
\textbf{17.} $g(x)=-(x+3)^2+5$; vertex $(-3,5)$

\medskip
\textbf{18.} (1) Vertical stretch $\times 4$\\
(2) Reflect over $x$-axis\\
(3) Shift left 1\\
(4) Shift down 3

\medskip
\textbf{19.} (1) Vertical shrink $\times\tfrac{1}{2}$\\
(2) Shift right 2\\
(3) Shift up 4

\medskip
\textbf{20.} $g(x)=-2(x-1)^2-4$

\medskip
\textbf{21.} Opens down ($a<0$); wider ($|a|=\tfrac{1}{2}$); shifted left 3 and up 5; vertex $(-3,5)$

\medskip
\textbf{22.} $g(x)=-3(x-2)^2+6$; vertex $(2,6)$

\medskip
\textbf{23.} $y=a(x+1)(x-5)$. Plug in $(2,-9)$:\\
$-9=a(3)(-3)=-9a \Rightarrow a=1$.\\
$y=(x+1)(x-5)=x^2-4x-5$

\medskip
\textbf{24.} $y=a(x-3)^2-2$. Plug in $(5,6)$:\\
$6=a(4)-2 \Rightarrow a=2$.\\
$y=2(x-3)^2-2=2x^2-12x+16$

\medskip
\textbf{25.} Vertex $x=\frac{-4+2}{2}=-1$.\\
$y=a(x+4)(x-2)$. Plug in $(-1,9)$:\\
$9=a(3)(-3)=-9a \Rightarrow a=-1$.\\
$y=-(x+4)(x-2)=-x^2-2x+8$

\medskip
\textbf{26.} $y=a(x-1)^2+4$. Plug in $(0,3)$:\\
$3=a(1)+4 \Rightarrow a=-1$.\\
$y=-(x-1)^2+4$\\
Domain: $(-\infty,+\infty)$; Range: $(-\infty,4]$

\medskip
\textbf{27a.} $h=-\frac{-6}{2}=3$, $k=9-18+5=-4$.\\
$y=(x-3)^2-4$

\medskip
\textbf{27b.} $h=-\frac{12}{6}=-2$, $k=12-24-1=-13$.\\
$y=3(x+2)^2-13$

\medskip
\textbf{27c.} $h=-\frac{4}{-2}=2$, $k=-4+8+12=16$.\\
$y=-(x-2)^2+16$

\medskip
\textbf{28.}
\begin{align*}
y &= 2(x^2-6x)+10\\
  &= 2(x^2-6x+9-9)+10\\
  &= 2(x-3)^2-18+10\\
  &= 2(x-3)^2-8
\end{align*}
Vertex $(3,-8)$

\medskip
\textbf{29a.} $f(x)=2(x-2)^2$; zero: $x=2$;\\
vertex $(2,0)$; domain $(-\infty,+\infty)$; range $[0,+\infty)$

\medskip
\textbf{29b.} Zeros $x=-5,1$; vertex $(-2,9)$;\\
domain $(-\infty,+\infty)$; range $(-\infty,9]$

\medskip
\textbf{29c.} Zeros $x=\pm\tfrac{1}{2}$; vertex $(0,-4)$;\\
domain $(-\infty,+\infty)$; range $[-4,+\infty)$

\medskip
\textbf{30.}
(a) $x = \frac{-3+7}{2} = 2$\\
(b) $y = a(x-2)^2+20$\\
(c) Use zero $x=7$: $0=a(25)+20 \Rightarrow a=-\tfrac{4}{5}$.\\
$y=-\tfrac{4}{5}(x-2)^2+20$\\
(d) Range: $(-\infty,\,20]$

\end{multicols}

\vspace{0.5cm}
\hrule
\begin{center}
  \textit{End of Quadratic Functions Notes \& Practice}
\end{center}

\end{document}
